\documentclass[../../main.tex]{subfiles}

\begin{document}
\rhead{東工大後期$2009$年$1$番}
\setcounter{equation}{0}
\setcounter{figure}{0}
\setcounter{table}{0}
\begin{oframed}
  \subfile{problem}
\end{oframed}
\setlength{\columnseprule}{0.4pt}
\begin{multicols}{2}
  {\bf [解答]}
  $P(X,Y,0)$, 球の中心$O' (0,0,1)$とする．
  球の様子は\cref{2009-1:fig:1}のようになる．

  \begin{figure}[H]
    \centering
    \begin{tikzpicture}
      \begin{axis}[
          axis lines=center,
          axis line style={gray!80!black, thick},
          xlabel=$x$,
          ylabel=$y$,
          zlabel=$z$,
          xlabel style={at={(axis cs:3.5,0,0)}, anchor=north east, font=\small},
          ylabel style={at={(axis cs:0,3.5,0)}, anchor=north west, font=\small},
          zlabel style={at={(axis cs:0,0,4.2)}, anchor=south, font=\small},
          xmin=-3, xmax=3.5,
          ymin=-2.5, ymax=3.5,
          zmin=0, zmax=4,
          xtick=\empty, ytick=\empty, ztick=\empty,
          view={125}{22},
          axis equal image,
          clip=false,
        ]

        % xy 平面上の影 S (楕円領域、斜線ハッチング)
        \addplot3[
          pattern=north east lines,
          pattern color=gray!60,
          draw=gray!80!black,
          thick,
          domain=0:360,
          samples=40,
        ] ( {-0.75 + 1.6*cos(x)}, {1.1*sin(x)}, {0} );
        \node[gray!80!black, font=\small] at (axis cs:-2.2, -0.8, 0) {$S$};

        % 球 (中心 (0,0,1), 半径 1)
        \addplot3[
          surf,
          domain=0:360, domain y=-90:90,
          opacity=0.35,
          samples=25, samples y=20,
          colormap/cool,
        ] ( {cos(y)*cos(x)}, {cos(y)*sin(x)}, {sin(y)+1} );

        % 球の中心 O'(0,0,1)
        \coordinate (Center) at (axis cs:0,0,1);
        \node[circle, fill=gray!80!black, inner sep=1.2pt] at (Center) {};
        \node[above left=2pt, font=\footnotesize] at (Center) {$O'(0,0,1)$};

        % 点光源 C(a,0,3) (a=1.5 の例)
        \coordinate (LightSource) at (axis cs:1.5,0,3);
        \node[circle, fill=red!85!black, inner sep=1.4pt] at (LightSource) {};
        \node[above right=2pt, red!85!black, font=\footnotesize] at (LightSource) {$C(a,0,3)$};

        % 影の点 P(X,Y,0)
        \coordinate (ShadowPoint) at (axis cs:-1.2, 0.4, 0);
        \node[circle, fill=blue!85!black, inner sep=1.4pt] at (ShadowPoint) {};
        \node[right=3pt, blue!85!black, font=\footnotesize] at (ShadowPoint) {$P(X,Y,0)$};

        % 線分 CP (光線)
        \draw[line width=1.1pt, blue!80!black] (LightSource) -- (ShadowPoint);

        % 光源 C から x 軸への垂線
        \draw[densely dotted, thick, gray!70!black] (LightSource) -- (axis cs:1.5,0,0) node[below=2pt, font=\footnotesize] {$a$};

        % 原点 O
        \node[below left=2pt, font=\small] at (axis cs:0,0,0) {$O$};

      \end{axis}
    \end{tikzpicture}
    \caption{球と点光源および影の領域$S$（$a=1.5$ の例）}
    \label{2009-1:fig:1}
  \end{figure}

  線分$CP$が球に遮られるためには，
  線分$CP$と球の中心$O'$の距離が球の半径である$1$より小さければ良い．
  直線$CP$はパラメータ表示で
  \begin{align}
    \mqty(x   \\y \\ z)
    = \mqty(a \\ 0 \\ 3) + t \mqty(X-a \\ Y \\ -3)\quad (t \in \mathbb{R})
  \end{align}
  だから，CP上の点$Q(t)$として，
  \begin{align}
    & |\overrightarrow{O'Q(t)}|^2                              \\
    = & \{ a+t(X-a) \}^2 + (tY)^2 + \{ 3-3t-1 \}^2    \\
    = & \{ a+t(X-a) \}^2 + (tY)^2 + (2-3t)^2          \\
    = & \{(X-a)^2+Y^2+9\}t^2 + 2\{a(X-a)-6\}t + a^2+4
  \end{align}
  となる．ここで簡単のため
  \begin{align}
    A & =(X-a)^2+Y^2+9(>0) \\
    B & =aX-a^2-6
  \end{align}
  とおいて，式変形を続けると
  \begin{align}
    |O'Q(t)|^2= & A\left(t+\dfrac{B}{A}\right)^2+a^2+4-\dfrac{B^2}{A}
  \end{align}
  だから，$t = - B/A$ の時最小値$a^2+4-B^2/A$をとる．
  よって条件は
  \begin{align}
    & \min |\overrightarrow{O'Q}|^2 \le 1                                    \\
    \Leftrightarrow & a^2+4 - \dfrac{\{a(X-a)-6\}^2}{(X-a)^2+Y^2+9} \le 1          \\
    \Leftrightarrow & a^2+3 - \dfrac{\{a(X-a)-6\}^2}{(X-a)^2+Y^2+9} \le 0          \\
    \Leftrightarrow & (a^2+3) \{(X-a)^2+Y^2+9\} - \{a(X-a)-6\}^2 \le 0            \\
    \Leftrightarrow & (a^2+3)(X-a)^2+(a^2+3)Y^2+9(a^2+3)                          \\
    & - \{a^2(X-a)^2-12a(X-a)+36\} \le 0   \\
    \Leftrightarrow & 3(X-a)^2+(a^2+3)Y^2+12a(X-a)+9a^2-9 \le 0                   \\
    \Leftrightarrow & 3\left[(X-a)+2a\right]^2+(a^2+3)Y^2-3(a^2+3) \le 0 \\
    \Leftrightarrow & \dfrac{(X+a)^2}{a^2+3} + \dfrac{Y^2}{3} \le 1
  \end{align}
  という楕円となり，これが求めるべき条件である．従って答えは
  \begin{align}
    \dfrac{(X+a)^2}{a^2+3} + \dfrac{Y^2}{3} \le 1
  \end{align}
  である．$\cdots$(答)

  {\bf [解説]}
  空間図形の問題．球体が点光源に対して作る影の形をとう問題で，我々は日常的な経験からこれが楕円になることを知っている．
  それを数学的に確かめなさいという面白い問題である．
  答案を見ればわかるように少し式変形が煩雑なので，答えが楕円になるという想定を持って変形を進めることで見通しが良くなるだろう．

  $CP$と$O'$の距離が$1$以下となる条件を立式する際，
  解答では直線$CP$上の点と$O'$の距離の最小値が$1$以下となるように立式した．
  ほぼ同じだが別の方法として，点と直線の距離を直接求めに行く方法も紹介しよう．
  一定検算に利用できるだろう．
  $O'$から直線$CP$におろした垂線の足を$H$とする．

  \begin{figure}[H]
    \centering
    \begin{tikzpicture}
      \begin{axis}[
          axis lines=none, % 軸を非表示
          xmin=0, xmax=10,
          ymin=0, ymax=8,
          ticks=none,
          clip=false,
        ]

        % --- 座標と直線の定義 ---
        % 点X1の座標
        \coordinate (X1) at (3,5);
        % 直線lが通る点x0
        \coordinate (x0) at (0,0);
        % 直線lの方向ベクトルa
        \coordinate (a) at (1,1);
        % 垂線の足X1'の座標を計算
        % X1' = x0 + t*a
        % ベクトル(X1 - X1') と a が直交する条件からtを求める
        % t = ((X1-x0)・a) / (a・a)
        \pgfmathsetmacro{\tvalue}{4}
        \coordinate (X1prime) at ($(x0) + \tvalue*(a)$);
        \coordinate (X2) at ($(x0) + 1.5*\tvalue*(a)$);

        % --- 描画 ---
        % 直線lを描画 (定義した2点x0とx0+aを通る長い線)
        \draw[name path=line] ($(x0) - 1*(a)$) -- ($(x0) + 8*(a)$);
        \draw[->] (x0) -- (X1);

        % 点とラベル
        \node[circle, fill, inner sep=1.5pt] at (x0) {}; \node[above=2pt] at (x0) {$C$};
        \node[circle, fill, inner sep=1.5pt] at (X1) {}; \node[above=2pt] at (X1) {$O'$};
        \node[circle, fill, inner sep=1.5pt] at (X1prime) {}; \node[below right=2pt] at (X1prime) {$H$};
        \node[circle, fill, inner sep=1.5pt] at (X2) {}; \node[above=2pt] at (X2) {$P$};

        % 点X1から直線lへの垂線
        \draw (X1) -- (X1prime);

        % 直角記号
        \coordinate (p1) at ($(X1prime)!4pt!(X1)$);
        \coordinate (p2) at ($(X1prime)!4pt!($(x0) + 1.5*(a)$)$);
        \coordinate (p3) at ($(p1)-(X1prime)+(p2)$);
        \draw (p1) -- (p3) -- (p2);

        % 距離dを示す弧とラベル
        \draw (X1) to [bend right=30] node[midway, left, xshift=-2pt] {$d$} (X1prime);

        % 直線のラベル
        \node[right=10pt] at ($(x0) + 1.5*(a)$) {$l_1$};

        % n1ベクトルを表す
        \coordinate (nstart) at (1,0);
        \coordinate (nend)   at (2,1);
        \draw[->] (nstart) -- (nend);
        \node[] at (1.5,0) {$\vec{n}$};

      \end{axis}
    \end{tikzpicture}
    \caption{直線$CP$と$O'$の距離．}
  \end{figure}

  $CH$はCO'のCPへの正射影ベクトルだから
  \begin{align}
    \overrightarrow{CH} = \dfrac{\overrightarrow{CP}\cdot\overrightarrow{CO'}}{|\overrightarrow{CP}|^2}\overrightarrow{CP}
  \end{align}
  であるから，
  \begin{align}
    \overrightarrow{O'H}
    & = \overrightarrow{CH}-\overrightarrow{CO'} \\
    & = \dfrac{\overrightarrow{CP}\cdot\overrightarrow{CO'}}{|\overrightarrow{CP}|^2}\overrightarrow{CP} - \overrightarrow{CO'}
  \end{align}
  となる．各ベクトルは
  \begin{align}
    \overrightarrow{CO'} & = \mqty(-a  \\ 0 \\ -2)\\
    \overrightarrow{CP}  & = \mqty(X-a \\ Y \\ -3)
  \end{align}
  だから，
  \begin{align}
    \overrightarrow{O'H}
    & = \dfrac{-a(X-a)+6}{(X-a)^2+Y^2+9}\mqty(X-a \\ Y \\ -3) - \mqty(-a  \\ 0 \\ -2) \\
    & = \dfrac{-B}{A}\mqty(X-a                    \\ Y \\ -3) - \mqty(-a  \\ 0 \\ -2)
  \end{align}
  この長さが$1$よりも小さければ良いので，
  \begin{align}
    \left|\overrightarrow{O'H}\right|^2 \le 1
  \end{align}
  が条件である．
  \begin{align}
    & \left|\overrightarrow{O'H}\right|^2                                                                          \\
    & = \left(\dfrac{-B}{A}(X-a)+a\right)^2+\left(\dfrac{-B}{A}Y\right)^2 + \left(\dfrac{3B}{A}+2\right)^2 \\
    & = \dfrac{1}{A^2}\left(-B(X-a)+aA\right)^2+\dfrac{B^2Y^2}{A^2} + \dfrac{1}{A^2}\left(3B+2A\right)^2
  \end{align}
  を得る．さらに整理すると
  \begin{align}
    \left|\overrightarrow{O'H}\right|^2
    & = \dfrac{B^2}{A^2}\left((X-a)^2+Y^2+9\right)                       \\
    & \quad \quad + \dfrac{-2aAB(X-a)}{A^2} + a^2 + 4 + \dfrac{12AB}{A^2} \\
    & = \dfrac{B^2}{A} + a^2 + 4 + \dfrac{-2B(a(X-a)-6)}{A}               \\
    & = \dfrac{B^2}{A} + a^2 + 4 + \dfrac{-2B^2}{A}                       \\
    & = -\dfrac{B^2}{A} + a^2 + 4
  \end{align}
  となり，解答の最小値と一致する．

  これでふた通りの方法でもとまった．
  導出を見ればわかるようにどちらも計算量には大差ないので，お好みで使い分けると良いだろう．

  \newpage
\end{multicols}
\end{document}
